\documentclass[11pt,a4paper]{article}
\usepackage[utf8]{inputenc}
\usepackage[T1]{fontenc}
\usepackage[spanish]{babel}
\usepackage[margin=2.2cm]{geometry}
\usepackage{amsmath,amssymb}
\usepackage{enumitem}
\usepackage{xcolor}
\usepackage{titlesec}
\usepackage{array}
\usepackage{booktabs}
\usepackage{longtable}
\usepackage{hyperref}

\hypersetup{colorlinks=true,urlcolor=blue,linkcolor=black}
\setlist{itemsep=2pt,topsep=2pt}
\titlespacing*{\section}{0pt}{10pt}{4pt}
\titlespacing*{\subsection}{0pt}{8pt}{3pt}

\title{\vspace{-1cm}\textbf{TP5 -- Sistema 1: Oscilador de Kuramoto}\\[2pt]\large Resumen teórico y modelo a implementar}
\author{Simulación de Sistemas -- 2026Q1}
\date{}

\begin{document}
\maketitle
\vspace{-0.8cm}

\section{Contexto neurocientífico}

El cerebro contiene $\sim 10^{11}$ neuronas con $\sim 10^{14}$--$10^{15}$ conexiones sinápticas. Cada neurona es una unidad de procesamiento con \textbf{dendritas} (entrada), \textbf{soma} (procesamiento no lineal) y \textbf{axón} (salida); se comunica con otras a través de \textbf{sinapsis} mediante \textbf{spikes} (pulsos de $\sim 100$ mV, $1$--$2$ ms). Lo central no es la forma del spike sino \emph{cuándo} ocurre \cite{ND-Ch1}.

Muchas neuronas son intrínsecamente oscilantes. Si idealizamos cada tren de spikes como una \textbf{fase} $\theta_i(t) \in [0,2\pi)$ que avanza a frecuencia natural $\omega_i$, podemos preguntar: ¿cuándo se sincroniza el colectivo?

\subsection*{Jerarquía de modelos neuronales}

\begin{enumerate}
  \item \textbf{Hodgkin--Huxley} (biofísico, $C_m \dot V = I - I_{Na} - I_K - I_L$): realista pero costoso.
  \item \textbf{FitzHugh--Nagumo} (reducido 2D, $\dot v = v - v^3/3 - w + I,\ \dot w = \epsilon(v+a-bw)$): captura excitabilidad y ciclos límite vía análisis del plano de fase \cite{ND-Ch4}.
  \item \textbf{Integrate-and-fire} ($\tau \dot V = -V + I$ con reset al umbral).
  \item \textbf{Kuramoto} (osciladores de fase acoplados): \textbf{este TP}.
\end{enumerate}

\section{Modelo de Kuramoto}

\subsection*{Ecuación de evolución}

Para $i = 1,\ldots,N$ con $N > 500$:
\begin{equation}
\frac{d\theta_i}{dt} \;=\; \omega_i \;+\; K \sum_{j} A_{ij}\, \sin(\theta_j - \theta_i)
\end{equation}

\begin{itemize}
  \item $\theta_i \in [0, 2\pi)$: fase de la neurona $i$.
  \item $\omega_i \sim \mathcal{N}(\mu=1,\ \sigma=0.1)$: frecuencia natural, \emph{fija} en cada realización.
  \item $K \in [0,1]$: acoplamiento global (parámetro de control).
  \item $A_{ij}$: matriz de conectividad (topología, $A_{ii}=0$).
  \item Condición inicial: $\theta_i(0) \sim \mathcal{U}[0, 2\pi)$.
\end{itemize}

\subsection*{Parámetro de orden}

Mide el grado de sincronización global:
\begin{equation}
r(t) \;=\; \left| \frac{1}{N} \sum_j \bigl[\cos\theta_j,\ \sin\theta_j\bigr] \right|
\;=\; \left| \frac{1}{N} \sum_j e^{i\theta_j} \right|
\end{equation}
$r \approx 0$: estado \emph{desincronizado}.\quad $r \approx 1$: \emph{sincronización global}.

\subsection*{Física: competencia heterogeneidad vs.\ acoplamiento}

El término $\sin(\theta_j - \theta_i)$ \textbf{atrae} fases cercanas, mientras que la dispersión de $\omega_i$ las \textbf{separa}. Existe un valor crítico $K_c$ a partir del cual emerge sincronización colectiva (transición de fase). Es análogo a la transición entre los regímenes \emph{Asynchronous Irregular} ($r\approx 0$) y \emph{Synchronous Regular} ($r\approx 1$) descritos en \cite{ND-Ch13} para redes de integrate-and-fire.

\subsection*{Topologías de red a estudiar}

\begin{description}[leftmargin=1.1cm,style=nextline]
  \item[(a) Red totalmente conectada] $A_{ij}=1\ \forall i\neq j$. \emph{All-to-all coupling}, el esquema más simple \cite{ND-Ch12}.
  \item[(b) Red aleatoria (Erdős--Rényi)] $A_{ij}=1$ con probabilidad $p\in[10^{-4},10^{-1}]$ (al menos 10 valores distribuidos logarítmicamente). Grado medio $\langle k \rangle = p(N-1)$.
  \item[(c) Red anillo] $A_{ij}=1$ sii $j \in \{i-v,\ldots,i-1,i+1,\ldots,i+v\}$ con índices periódicos, $v\in[1,10]$.
\end{description}

\section{Parámetros del Sistema 1 (según enunciado)}

\subsection*{Valores numéricos elegidos para este TP}

Calibrados experimentalmente (ver \texttt{calibracion.pdf} para todo el detalle):

\begin{center}
\renewcommand{\arraystretch}{1.2}
\begin{tabular}{@{}l l p{8cm}@{}}
\toprule
\textbf{Parámetro} & \textbf{Valor} & \textbf{Justificación} \\
\midrule
$N$ & \textbf{600} & Cumple $N>500$, costo $\sim 8\times$ menor que $N=1000$. \\
$dt$ & \textbf{$10^{-3}$} & Mínimo $dt = 10^{-x}$ estable de RK4, precisión de máquina. \\
$t_{\text{sim}}$ completa  & \textbf{50}   & Sincroniza en $t<1$. \\
$t_{\text{sim}}$ aleatoria & \textbf{50}   & Sincroniza en $t<3$ incluso para $p$ chico. \\
$t_{\text{sim}}$ anillo    & \textbf{1500} & Transitorios largos por propagación local (estados quiméricos). \\
\bottomrule
\end{tabular}
\end{center}

\subsection*{Parámetros comunes a las tres topologías}

\begin{center}
\renewcommand{\arraystretch}{1.25}
\begin{tabular}{@{}p{1.4cm} p{4.6cm} p{4.6cm} p{4.6cm}@{}}
\toprule
\textbf{Símbolo} & \textbf{Significado} & \textbf{Valor / Rango} & \textbf{Cómo se determina} \\
\midrule
$N$ & Número de osciladores & $N = 600$ (consigna: $N>500$) & Calibrado por costo vs.\ calidad estadística. \\
$\omega_i$ & Frecuencia natural & $\omega_i \sim \mathcal{N}(\mu{=}1,\ \sigma{=}0.1)$ & Muestreo aleatorio por realización; \emph{fijas} durante la integración. \\
$\theta_i(0)$ & Fase inicial & $\theta_i(0) \sim \mathcal{U}[0, 2\pi)$ & Muestreo aleatorio uniforme por realización. \\
$K$ & Acoplamiento global & $K \in [0, 1]$ & Parámetro de control; se \emph{barre}. \\
$A_{ij}$ & Conectividad & Depende de la topología & $A_{ii}=0$ siempre. \\
$dt$ & Paso de integración & $dt = 10^{-3}$ (fijo) & Calibrado por test de convergencia RK4. \\
$t_{\text{sim}}$ & Tiempo total simulado & Por topología (ver arriba) & Calibrado por análisis de estacionariedad de $r(t)$. \\
\bottomrule
\end{tabular}
\end{center}

\subsection*{Parámetros específicos por topología}

\textbf{(a) Red totalmente conectada:}
$A_{ij}=1\ \forall i\neq j$. Barrer $K\in[0,1]$. Para cada $K$: \textbf{$>10$ realizaciones} de $\theta_i(0)$ y $\omega_i$, promediar $r(t)$.

\medskip
\textbf{(b) Red aleatoria (Erdős--Rényi):}
$A_{ij}=1$ con probabilidad $p$, $0$ en otro caso ($\forall i\neq j$).
\begin{itemize}
  \item $p \in [10^{-4}, 10^{-1}]$, \textbf{al menos 10 valores distribuidos logarítmicamente}.
  \item Para análisis de $r(p)$: fijar $K=0.1$.
  \item Por cada par $(p, K)$: \textbf{$>10$ realizaciones} independientes de la red $A_{ij}$ \emph{y} de las condiciones iniciales.
  \item Mapas 2D: $r_\infty$ y $\tau_s$ como función de $(p, K)$.
\end{itemize}

\medskip
\textbf{(c) Red anillo:}
$A_{ij}=1$ sii $j \in \{i-v,\ldots,i-1,i+1,\ldots,i+v\}$, índices periódicos (mod $N$).
\begin{itemize}
  \item $v \in [1, 10]$ entero (vecindad a cada lado).
  \item Para análisis de $r(v)$: fijar $K=0.1$.
  \item Por cada par $(v, K)$: \textbf{$>10$ realizaciones} independientes de las condiciones iniciales.
  \item Mapas 2D: $r_\infty$ y $\tau_s$ como función de $(v, K)$.
\end{itemize}

\subsection*{Observación: simetría de $A_{ij}$}

El enunciado no la fija explícitamente. Convención adoptada:
\begin{itemize}
  \item \emph{Completa}: simétrica trivialmente.
  \item \emph{Aleatoria}: sortear $A_{ij}$ para $i<j$ y replicar $A_{ji}=A_{ij}$ (red no dirigida).
  \item \emph{Anillo}: simétrica por construcción.
\end{itemize}

\section{Integrador: Runge--Kutta de orden 4}

Con $\boldsymbol\theta = (\theta_1,\ldots,\theta_N)$ y campo $f(\boldsymbol\theta)_i = \omega_i + K \sum_j A_{ij}\sin(\theta_j-\theta_i)$:

\begin{align*}
\mathbf k_1 &= f(\boldsymbol\theta) \\
\mathbf k_2 &= f(\boldsymbol\theta + \tfrac{h}{2}\,\mathbf k_1) \\
\mathbf k_3 &= f(\boldsymbol\theta + \tfrac{h}{2}\,\mathbf k_2) \\
\mathbf k_4 &= f(\boldsymbol\theta + h\,\mathbf k_3) \\
\boldsymbol\theta(t+h) &= \boldsymbol\theta(t) + \tfrac{h}{6}\bigl(\mathbf k_1 + 2\mathbf k_2 + 2\mathbf k_3 + \mathbf k_4\bigr)
\end{align*}

El sistema es \emph{autónomo} (f no depende explícitamente de $t$), por lo que los offsets temporales se omiten. $h = dt$ fijo durante toda la simulación.

\subsection*{Resultado de la calibración (ver \texttt{calibracion.pdf})}

\textbf{$dt$:} test de convergencia con $dt = 10^{-x}$ sobre red completa $N=1000$, $K=0.5$. Para $dt \geq 10^{-2}$ el integrador RK4 es \emph{inestable} (error $\sim 1$); para $dt \leq 10^{-3}$ converge a precisión de máquina ($\sim 10^{-15}$). $\Rightarrow dt = 10^{-3}$.

\textbf{$t_{\text{sim}}$ por topología} (necesario por la física, no por capricho):
\begin{itemize}
\item \emph{Completa} y \emph{aleatoria}: sincronizan en $t<3$ incluso para $p$ chico. $t_{\text{sim}} = 50$ deja margen $>15\times$.
\item \emph{Anillo}: propagación local $\Rightarrow \tau \sim N/v$, plus estados quiméricos. Confirmamos con corridas a $t_{\text{sim}}=2000$ que $r(t)$ recién se estabiliza alrededor de $t\approx 1000$--$1500$ para $v$ intermedio. $t_{\text{sim}} = 1500$ captura el estacionario.
\end{itemize}

\section{Observables y objetivos}

\begin{enumerate}
  \item Evolución temporal $r(t)$ para distintos $K\in[0,1]$, promediada sobre $>10$ realizaciones independientes (semillas de $\theta_i(0)$ y, si corresponde, de $A_{ij}$).
  \item Curva estacionaria $r_\infty(K)$: identificar $K_c$.
  \item Tiempo de sincronización $\tau_s(K)$ (llegada al estacionario).
  \item Mapas 2D: $r_\infty(p,K)$ y $\tau_s(p,K)$ para red aleatoria; $r_\infty(v,K)$ y $\tau_s(v,K)$ para anillo.
  \item Comparación final de $\tau_s$ entre las tres topologías.
\end{enumerate}

\section{Plan de implementación}

\begin{enumerate}
  \item Motor de simulación: integra con RK4, $dt$ fijo, emite \textbf{archivo de texto} (consigna).
  \item Análisis y animación: módulos separados que leen los archivos de texto.
  \item Por realización: muestrear $\omega_i$, $\theta_i(0)$ y (si corresponde) $A_{ij}$; integrar hasta tiempo final; loggear $r(t)$ y/o todas las $\theta_i(t)$.
  \item Promediar sobre $\geq 10$ realizaciones por punto de $(K,p)$ o $(K,v)$.
  \item Visualización: nodos en el círculo unitario coloreados por fase, o grilla coloreada por $\theta_i$.
\end{enumerate}

\section*{Bibliografía}
\begin{thebibliography}{9}
\bibitem{ND-Ch1} W.~Gerstner, W.~Kistler, R.~Naud, L.~Paninski. \emph{Neuronal Dynamics}. Cambridge University Press. Cap.~1: \emph{Introduction: Neurons and Mathematics}. \url{https://neuronaldynamics.epfl.ch/online/Ch1.html}
\bibitem{ND-Ch4} \emph{ibid.} Cap.~4: \emph{Dimensionality Reduction and Phase Plane Analysis}. \url{https://neuronaldynamics.epfl.ch/online/Ch4.html}
\bibitem{ND-Ch12} \emph{ibid.} Cap.~12.3: \emph{Connectivity Schemes}. \url{https://neuronaldynamics.epfl.ch/online/Ch12.S3.html}
\bibitem{ND-Ch13} \emph{ibid.} Cap.~13.4: \emph{Synchrony, oscillations, and irregularity}. \url{https://neuronaldynamics.epfl.ch/online/Ch13.S4.html}
\bibitem{T5a} Cátedra SdS. \emph{Teorica5a: Dinámica neuronal y sincronización}. (Slides 1--14.)
\bibitem{E} Cátedra SdS. \emph{TP5: Comportamiento colectivo -- Enunciado}, 22/05/2026.
\end{thebibliography}

\end{document}
